\textbf{Phần 1. Trắc nghiệm lựa chọn}
\setcounterpageref{ex}{0}

\begin{ex} %[3D3V67]
Bạn Mai rời nhà lúc 6 giờ 45 phút và đến trường lúc 7 giờ 10 phút. Thời gian Mai đi từ nhà đến trường là
\choice
{$15$ phút}
{$35$ phút}
{\True $25$ phút}
{$20$ phút}
\loigiai{
Từ 6h45 đến 7h00 là 15 phút, từ 7h00 đến 7h10 là 10 phút. Tổng thời gian: $15 + 10 = 25$ phút.
\textbf{Đáp án của câu hỏi là C}
}
\end{ex}

\begin{ex} %[3D3V67]
Nếu hôm nay là thứ Tư ngày 15 thì thứ Tư tuần sau là ngày nào?
\choice
{\True Ngày $22$}
{Ngày $21$}
{Ngày $23$}
{Ngày $20$}
\loigiai{
Thứ Tư tuần sau là: $15 + 7 = 22$.
\textbf{Đáp án của câu hỏi là A}
}
\end{ex}

\begin{ex} %[3D3V67]
Đồng hồ điện tử hiển thị $18:45$. Thời gian này tương ứng với
\choice
{$6$ giờ $45$ phút sáng}
{$8$ giờ kém $15$ phút tối}
{$18$ giờ $45$ phút sáng}
{\True $7$ giờ kém $15$ phút tối}
\loigiai{
$18$ giờ là 6 giờ chiều; $18:45$ là 6 giờ 45 phút chiều hay 7 giờ kém 15 phút tối.
\textbf{Đáp án của câu hỏi là D}
}
\end{ex}

\begin{ex} %[3D3V67]
Một buổi học bắt đầu từ 8 giờ sáng và kéo dài trong 3 tiết học, mỗi tiết 40 phút, giờ ra chơi giữa các tiết tổng cộng 20 phút. Buổi học kết thúc lúc
\choice
{$10$ giờ}
{\True $10$ giờ $20$ phút}
{$10$ giờ $10$ phút}
{$10$ giờ $40$ phút}
\loigiai{
Tổng thời gian học và chơi: $40 \times 3 + 20 = 140$ phút = 2 giờ 20 phút. Kết thúc: $8\text{h} + 2\text{h } 20\text{p} = 10$ giờ 20 phút.
\textbf{Đáp án của câu hỏi là B}
}
\end{ex}

\begin{ex} %[3D3V67]
Ngày Quốc tế Thiếu nhi 1 tháng 6 năm nay là thứ Tư. Hỏi ngày 8 tháng 6 cùng năm là thứ mấy?
\choice
{Thứ Ba}
{Thứ Năm}
{\True Thứ Tư}
{Chủ nhật}
\loigiai{
Cách nhau đúng 7 ngày ($8 - 1 = 7$) nên cũng là ngày thứ Tư.
\textbf{Đáp án của câu hỏi là C}
}
\end{ex}

\begin{ex} %[3D3V67]
Tàu hỏa chạy từ ga Hà Nội lúc 22 giờ đêm hôm trước và đến ga Vinh lúc 4 giờ 30 phút sáng hôm sau. Thời gian tàu chạy là
\choice
{\True $6$ giờ $30$ phút}
{$5$ giờ $30$ phút}
{$6$ giờ}
{$7$ giờ}
\loigiai{
Từ 22h đến 24h là 2 giờ; từ 0h đến 4h30 là 4 giờ 30 phút. Tổng thời gian: $2\text{h} + 4\text{h } 30\text{p} = 6$ giờ 30 phút.
\textbf{Đáp án của câu hỏi là A}
}
\end{ex}

\begin{ex} %[3D3V67]
Một năm có bao nhiêu quý, mỗi quý có mấy tháng?
\choice
{$3$ quý, mỗi quý có $4$ tháng}
{$2$ quý, mỗi quý có $6$ tháng}
{$4$ quý, mỗi quý có $4$ tháng}
{\True $4$ quý, mỗi quý có $3$ tháng}
\loigiai{
Một năm có 4 quý, mỗi quý có $12 : 4 = 3$ tháng.
\textbf{Đáp án của câu hỏi là D}
}
\end{ex}

\begin{ex} %[3D3V67]
Kim phút của đồng hồ quay được đúng 2 vòng rưỡi thì kim giờ quay được bao nhiêu phần của một vòng đồng hồ?
\choice
{5 khoảng giờ}
{\True Hai khoảng giờ rưỡi (tương ứng $2$ giờ $30$ phút)}
{1 vòng}
{Nửa vòng}
\loigiai{
Kim phút quay 1 vòng là 1 giờ. Quay 2 vòng rưỡi tức là 2 giờ 30 phút (kim giờ đi được 2 khoảng rưỡi).
\textbf{Đáp án của câu hỏi là B}
}
\end{ex}

\begin{ex} %[3D3V67]
Ngày cuối cùng của tháng 7 luôn luôn là ngày
\choice
{Ngày $30$}
{Ngày $28$}
{\True Ngày $31$}
{Ngày $29$}
\loigiai{
Tháng 7 có 31 ngày nên ngày cuối cùng là ngày 31.
\textbf{Đáp án của câu hỏi là C}
}
\end{ex}

\begin{ex} %[3D3V67]
Một chiếc đồng hồ báo thức reo chuông cách mỗi 15 phút một lần. Lần đầu reo lúc 6 giờ đúng. Lần thứ tư reo lúc mấy giờ?
\choice
{\True $6$ giờ $45$ phút}
{$7$ giờ}
{$6$ giờ $30$ phút}
{$7$ giờ $15$ phút}
\loigiai{
Lần 1: 6h00; Lần 2: 6h15; Lần 3: 6h30; Lần 4: 6h45.
\textbf{Đáp án của câu hỏi là A}
}
\end{ex}

\begin{ex} %[3D3V67]
Tháng 2 của một năm có 5 ngày Chủ nhật. Ngày Chủ nhật đầu tiên của tháng 2 đó chắc chắn là ngày
\choice
{Ngày $2$}
{Ngày $3$}
{Ngày $7$}
{\True Ngày $1$}
\loigiai{
Để có 5 ngày Chủ nhật thì tháng 2 phải có 29 ngày (năm nhuận) và ngày đầu tiên (ngày 1) phải là Chủ nhật (các CN là ngày 1, 8, 15, 22, 29).
\textbf{Đáp án của câu hỏi là D}
}
\end{ex}

\begin{ex} %[3D3V67]
Một người thợ bắt đầu làm việc lúc 7 giờ 30 phút, nghỉ trưa 1 giờ 30 phút và kết thúc công việc lúc 17 giờ. Thời gian làm việc thực tế của người thợ là
\choice
{$9$ giờ $30$ phút}
{\True $8$ giờ}
{$8$ giờ $30$ phút}
{$7$ giờ $30$ phút}
\loigiai{
Tổng thời gian từ 7h30 đến 17h00 là 9 giờ 30 phút. Trừ 1 giờ 30 phút nghỉ trưa: $9\text{h } 30\text{p} - 1\text{h } 30\text{p} = 8$ giờ.
\textbf{Đáp án của câu hỏi là B}
}
\end{ex}

\textbf{Phần 2. Trắc nghiệm đúng sai}
\setcounterpageref{ex}{0}

\begin{ex} %[3D3V67]
Xét tính đúng sai của việc thực hành xem lịch và đồng hồ:
\choiceTF
{\True Kim dài của đồng hồ quay 1 vòng thì hết đúng 60 phút (1 giờ).}
{Một tuần lễ có 5 ngày làm việc và không có ngày nào khác.}
{\True Hai ngày Chủ nhật liền kề nhau luôn cách nhau đúng 7 ngày.}
{Lúc 12 giờ trưa hai kim đồng hồ vuông góc với nhau.}
\loigiai{
\begin{itemize}
    \item Ý a) Đúng: Đặc tính hoạt động của đồng hồ.
    \item Ý b) Sai: Sai vì tuần lễ có 7 ngày.
    \item Ý c) Đúng: Quy luật tuần lễ.
    \item Ý d) Sai: Sai vì lúc 12 giờ hai kim trùng nhau.
\end{itemize}
\begin{itemize}
    \item a) \textbf{Đúng}.
    \item b) \textbf{Sai}.
    \item c) \textbf{Đúng}.
    \item d) \textbf{Sai}.
\end{itemize}
}
\end{ex}

\begin{ex} %[3D3V67]
Một chuyến du lịch kéo dài từ sáng ngày 28 tháng 5 đến hết chiều ngày 2 tháng 6 cùng năm:
\choiceTF
{\True Tháng 5 có 31 ngày.}
{\True Trong tháng 5 đoàn đi du lịch trong 4 ngày (28, 29, 30, 31).}
{\True Trong tháng 6 đoàn đi du lịch trong 2 ngày (1, 2).}
{Chuyến du lịch kéo dài đúng 10 ngày.}
\loigiai{
\begin{itemize}
    \item Ý a) Đúng: Kiến thức chuẩn.
    \item Ý b) Đúng: Bốn ngày cuối tháng 5.
    \item Ý c) Đúng: Hai ngày đầu tháng 6.
    \item Ý d) Sai: Sai vì tổng cộng là $4 + 2 = 6$ ngày.
\end{itemize}
\begin{itemize}
    \item a) \textbf{Đúng}.
    \item b) \textbf{Đúng}.
    \item c) \textbf{Đúng}.
    \item d) \textbf{Sai}.
\end{itemize}
}
\end{ex}

\textbf{Phần 3. Trả lời ngắn}
\setcounterpageref{ex}{0}

\begin{ex} %[3D3V67]
Một trận đấu cờ vua bắt đầu lúc 14 giờ 15 phút và kết thúc lúc 15 giờ 05 phút. Trận đấu diễn ra trong bao nhiêu phút?
\shortans{50}
\loigiai{
\textbf{Đáp số:} $50$\\
Từ 14h15 đến 15h05 là 50 phút.
}
\end{ex}

\begin{ex} %[3D3V67]
Tháng 10 có bao nhiêu ngày?
\shortans{31}
\loigiai{
\textbf{Đáp số:} $31$\\
Tháng 10 có 31 ngày.
}
\end{ex}

\begin{ex} %[3D3V67]
Nếu ngày 20 tháng 11 là Chủ nhật thì ngày 27 tháng 11 cùng năm là ngày thứ mấy (Chủ nhật quy ước ghi 8)?
\shortans{8}
\loigiai{
\textbf{Đáp số:} $8$\\
$20 + 7 = 27$ nên vẫn là Chủ nhật (8).
}
\end{ex}

\begin{ex} %[3D3V67]
Đổi 3 giờ ra đơn vị phút.
\shortans{180}
\loigiai{
\textbf{Đáp số:} $180$\\
$3 \times 60 = 180$ phút.
}
\end{ex}

\begin{ex} %[3D3V67]
Một chương trình truyền hình kéo dài từ 20 giờ 00 phút đến 20 giờ 45 phút. Chương trình kéo dài bao nhiêu phút?
\shortans{45}
\loigiai{
\textbf{Đáp số:} $45$\\
45 phút.
}
\end{ex}

\begin{ex} %[3D3V67]
Trong một năm có bao nhiêu tháng có 31 ngày?
\shortans{7}
\loigiai{
\textbf{Đáp số:} $7$\\
Gồm 7 tháng: 1, 3, 5, 7, 8, 10, 12.
}
\end{ex}
